\section{Electron counts in large nuclear voids}\label{m:sec:void}

The local count bound of Lemma~\ref{m:lem:screen-count} controls the
electrons in each annulus whose radius is comparable to its distance
from the nuclei.  Summing that bound over arbitrarily many annuli would
lose a factor equal to the number of scales.  We first remove this loss
in a nuclear void.  A geometric decomposition will then reduce the
exterior region of an arbitrary molecule to only $O(M)$ such voids.

Throughout this section $\Psi$ is the normalized eigenstate in the
largest strictly bound sector, and $E=E_N=\inf_{j\ge0}E_j$, as in
Lemma~\ref{lem:sector-global-minimum}.  Expectations without an additional
subscript refer to its position law, with spins summed.  Constants may
depend on the fixed numbers $A=40$ and $L=10^6$, but on no feature of the
molecule.

\begin{lemma}[A large nuclear void]\label{m:lem:void-shell}
There are universal constants $t_0,C>0$ with the following property.
Let $O\in\R^3$, $t\ge t_0$ and $T=2^Jt$, where $J\ge1$ is an integer.
Suppose that at least one nucleus lies in $B(O,t/100)$ and that every
nucleus satisfies
\[
 |R_m-O|<t/100\qquad\text{or}\qquad |R_m-O|>100T.
\]
Then
\begin{equation}\label{m:eq:void-annulus-count}
 \left\|\#\{i:t\le |x_i-O|<T\}\right\|_{L^2(\Psi)}\le C.
\end{equation}
\end{lemma}

\begin{proof}
Translate $O$ to the origin.  Put $u_i=|x_i|$, $s_k=2^kt$ and
\[
 n_k=\#\{i:s_k\le u_i<2s_k\},\qquad
 L_{\mathrm{ann}}=\sum_{k=0}^{J-1}n_k.
\]
For $s_k/4\le |v|\le8s_k$, the nuclear assumptions give
$c s_k\le d(v)\le C s_k$, where $d(v)$ is the distance to the nearest
nucleus.  Indeed, the inner nucleus supplies the upper bound; the inner
and outer alternatives supply the lower bound.  Choose $t_0$ so large
that $d(v)/L>1$ on all these annuli.  Lemma~\ref{m:lem:screen-scale} then
gives $a_v=d(v)/L\asymp s_k/L$.

Each enlarged annulus is covered by a universally bounded number of
local balls $B(v,a_v)$ with comparable scales.  For example, take a
maximal set of centers in the annulus separated by $c s_k/L$, with $c$
small; the disjoint balls of half that radius give the bound on their
number.  Lemma~\ref{m:lem:screen-count}, applied to any $\delta$-state
$\omega$ in the fixed sector $N$, consequently yields
\begin{equation}\label{m:eq:void-local-shell}
 \|n_k\|_{L^2(\omega)}
 +\|n_{\mathrm{neigh},k}\|_{L^2(\omega)}
 \le C(1+\sqrt{\delta s_k}),
 \qquad
 n_{\mathrm{neigh},k}=\#\{i:s_k/2\le u_i\le4s_k\}.
\end{equation}
This proves the result when $J$ is bounded.  We may therefore assume that
$J$ exceeds a fixed sufficiently large integer.

\medskip\noindent
\emph{A weighted eigenstate inequality.}
Choose $0\le\chi\le1$ supported in $(32t,T/32)$, equal to one on
$[64t,T/64]$, and satisfying
$|\chi'(u)|\le C/u$, $|\chi''(u)|\le C/u^2$.
Such a function is obtained by multiplying a fixed cutoff at scale $t$
by a fixed cutoff at scale $T$.  Set $w(u)=u\chi(u)^2$.
It is a bounded smooth multiplier that vanishes near the origin.
For each weighted index, integration by parts gives
\begin{align*}
 \operatorname{Re}\langle w(u_i)\Psi,-\tfrac12\Delta_i\Psi\rangle
 &=\tfrac12\int w(u_i)|\nabla_i\Psi|^2
   -\tfrac14\int \Delta_iw(u_i)|\Psi|^2\\
 &\ge-C\E\frac{\1_{\{t\le u_i<T\}}}{u_i}.
\end{align*}
These identities hold first after form-domain approximation and then
by continuity, since $w$ and its derivatives are bounded for fixed
$t,T$.  Summing the last bound costs at most
$C\sum_k s_k^{-1}\E n_k\le C/t\le C$ by
\eqref{m:eq:void-local-shell} with $\delta=0$.
After fixing $x_i$ and its spin, the remaining Hamiltonian has energy at
least $E_{N-1}\ge E$.  Its contribution against the nonnegative weight
$w(u_i)$ is therefore nonnegative.  Testing the eigen-equation with
$\sum_iw(u_i)\Psi$ shows that weighted repulsion minus attraction is at
most $C$.

For an electron in a shell meeting the support of $\chi$, keep only its
repulsion against electrons at radius $<t$, at radius $\ge T$, and in
shells $l\le k-5$.  Dropping the other repulsions preserves the upper
bound.  Write
\[
 P_k=\sum_{\substack{0\le l<J\\l\le k-5}}n_l,
 \qquad q=\nu(B(0,t))-\#\{i:u_i<t\},
\]
and define the exterior signed potential at the origin by
\[
 B=\int_{|R|>100T}\frac{d\nu(R)}{|R|}
       -\sum_{u_i\ge T}\frac1{u_i}.
\]
For each such $k$, use a cut that is full out on
$[s_k/2,4s_k]$, supported in $[s_k/4,8s_k]$, and has gradients bounded
by $C/s_k$.  Let $X_k$ be its out data as in
Lemma~\ref{lem:loc-identity}; a bar will mean conditional expectation
given $X_k$ in this paragraph and below.  Set
\[
 \varphi(v)=V(v)-\sum_{u_i<t}\frac1{|v-x_i|}
                   -\sum_{u_i\ge T}\frac1{|v-x_i|},
 \qquad
 \psi_k(v)=\overline{\varphi(v)
            -\sum_{l\le k-5}\sum_{i\in l}\frac1{|v-x_i|}}.
\]
Every supported shell satisfies $32t\le s_k\le T/64$.
All sources defining $\psi_k$ are outside
$\{s_k/4<|v|<8s_k\}$: prefix electrons have radius $<s_k/16$,
inner electrons have radius $<t\le s_k/32$, and exterior electrons have
radius $\ge T\ge64s_k$.  Thus $\psi_k$ has a continuous harmonic version
on this annulus.  Newton's sphere-average formula gives
\begin{equation}\label{m:eq:void-sphere-mean}
 \langle\psi_k\rangle_u
   =\frac{\bar q-\bar P_k}{u}+\bar B,
 \qquad s_k/4<u<8s_k,
\end{equation}
where $\langle F\rangle_u=(4\pi u^2)^{-1}\int_{|v|=u}F(v)\,dS(v)$.
The shell-$k$ positions are all recorded in $X_k$.  We may therefore
condition the selected interaction terms while retaining their weights.
The preceding eigenstate inequality becomes
\begin{equation}\label{m:eq:void-weighted-field}
 -\E\sum_k\sum_{i\in k}u_i\chi(u_i)^2\psi_k(x_i)\le C,
\end{equation}
where $k$ ranges over the supported shells.  Conditional harmonic
versions and sphere averages are justified directly by the positive
distance to these sources and their fixed finite total masses.

\medskip\noindent
\emph{A size-biased law for one shell.}
Fix a supported $k$.  Choose a smooth function $b_k$, equal to one on
$[s_k,2s_k]$, supported in $(s_k/2,4s_k)$, with
$0\le b_k\le1$ and $|b_k'|\le C/s_k$.  Put
\[
 S_k=\sum_i b_k(u_i)^2,\qquad m_k=\E S_k.
\]
Then $n_k\le S_k\le n_{\mathrm{neigh},k}$.  If $m_k=0$, every
shell-$k$ contribution vanishes and can be omitted.  Otherwise define
\[
 \Psi_k=\sqrt{S_k/m_k}\,\Psi,
\]
and denote its expectation and $L^2$ norm by $\E_k$ and $\|\cdot\|_{2,k}$.
The square root of a sum of squares is Lipschitz, and, off its zero set,
\[
 |\nabla\sqrt{S_k}|^2
 =\frac{\sum_i b_k(u_i)^2|\nabla b_k(u_i)|^2}{S_k}
 \le Cs_k^{-2}\1_{\{n_{\mathrm{neigh},k}>0\}}
 \le Cs_k^{-2}n_{\mathrm{neigh},k}.
\]
Its gradient is zero almost everywhere on its zero set.  The multiplier
identity in Lemma~\ref{lem:sector-multiplier} shows that $\Psi_k$ is a
$\delta_k$-state, with
\begin{equation}\label{m:eq:void-tilt-energy}
 \delta_k\le
 \frac{C\E n_{\mathrm{neigh},k}}{s_k^2m_k}.
\end{equation}
The random variable $S_k$ is $X_k$-measurable, because its support is
inside the full-cut zone.  The joint position and cut-label law changes
by the factor $S_k/m_k$.  Its conditional position law given $X_k$ is
therefore unchanged on $S_k>0$.  In particular all the barred fields and
counts above have the same conditional versions under this new law.
This permits us to use the local estimates for $\Psi_k$ without changing
the fields in \eqref{m:eq:void-weighted-field}.

\medskip\noindent
\emph{A positive bound and a signed harmonic comparison.}
On the broad annulus of harmonicity, the base and prefix sources lie
outside the exclusion defining $J_v$.  Its radius $Aa_v$ is a fixed
small fraction of $|v|$, since $A/L$ is small.  The additional electrons
subtracted by $J_v$ give the upper bound
\begin{equation}\label{m:eq:void-positive-comparison}
 \psi_k(v)\le\overline{J_v}
       +C\sum_{l\ge k-4}\frac{\bar n_l}{\max(s_l,s_k)}.
\end{equation}
For comparable shells the exclusion supplies the denominator
$c s_k$; for sufficiently larger shells radial separation supplies
$c s_l$.  These are precisely the electrons absent from the definition
of $\psi_k$.

The cut defining $X_k$ is covered by a universally bounded number of
cells at scales comparable to $s_k/L$.  Its gradients are $C/s_k$ and
$a_v\ge1$ there, so it meets the initial-cut hypothesis of
Lemma~\ref{m:lem:screen-positive-field}.  For targets in the same broad
annulus, that lemma and \eqref{m:eq:void-local-shell}, now applied under
$\Psi_k$, give
\[
 \|(\psi_k(v))_+\|_{2,k}
 \le \frac{C(1+\sqrt{\delta_k s_k})}{s_k}.
\]
Conditional Jensen's inequality bounds $\|\bar n_l\|_{2,k}$ by
$\|n_l\|_{2,k}$.  Thus the added count terms are bounded by the convergent series
\[
 C\sum_{l\ge k-4}\frac{1+\sqrt{\delta_k s_l}}{s_l}
 \le \frac{C(1+\sqrt{\delta_k s_k})}{s_k}.
\]
Positive parts of harmonic functions are subharmonic.  Covering the
compact annulus $\{s_k/2\le|v|\le3s_k\}$ by a fixed number of interior
balls in the broad annulus, and using the ball submean inequality and
Minkowski's inequality, consequently gives
\begin{equation}\label{m:eq:void-positive-supremum}
 A_k=s_k\sup_{s_k/2\le|v|\le3s_k}(\psi_k(v))_+,
 \qquad \|A_k\|_{2,k}\le C(1+\sqrt{\delta_k s_k}).
\end{equation}
The function $A_k/s_k-\psi_k$ is nonnegative and harmonic in the interior
of this latter annulus.  Its value at any point of a sphere
$s_k\le u\le2s_k$ is at least $c$ times its sphere mean, for one universal
$c>0$.  Indeed, after scaling, a bounded chain of interior balls joins
any two points of these spheres; on a ball, this comparison follows by
comparing its Poisson kernels on a smaller concentric ball.  Applying
this comparison and \eqref{m:eq:void-sphere-mean} gives
\begin{equation}\label{m:eq:void-harnack}
 u\psi_k(v)\le C A_k+c(\bar q+u\bar B-\bar P_k),
 \qquad |v|=u\in[s_k,2s_k].
\end{equation}
The negative prefix term retains a fixed positive coefficient.

\medskip\noindent
\emph{The endpoint fields.}
Put $s_-=2t$ and $s_+=T/2$.  For a supported shell and $u\in[s_k,2s_k]$,
\[
 \bar q+u\bar B
 =(1-\vartheta(u))(\bar q+s_-\bar B)
      +\vartheta(u)(\bar q+s_+\bar B),
 \quad
 \vartheta(u)=\frac{u-s_-}{s_+-s_-},
 \quad 0\le\vartheta(u)\le C s_k/T.
\]
At either endpoint all base sources are outside the $J_v$ exclusion.
Sphere averaging $\bar\varphi$ and adding back the annular electrons
therefore yields
\begin{equation}\label{m:eq:void-endpoint-field}
 \bar q+s\bar B\le
 s\langle(\overline{J_v})_+\rangle_s
       +Cs\sum_l\frac{\bar n_l}{\max(s_l,s)},
 \qquad s\in\{s_-,s_+\}.
\end{equation}
Indeed, the omitted contribution of an annular electron is nonnegative
and at most its full Coulomb kernel.  The latter has sphere average
$1/\max(s,u_i)\le1/\max(s,s_l)$ for an electron in shell $l$.
For these targets $a_v\asymp s/L$.  The initial cut is still at scale
$s_k/L$, so the energy parameter in
Lemma~\ref{m:lem:screen-positive-field} satisfies
$e_*\le C(\min(s,s_k)^{-1}+\delta_k)$.  The precise consequence is
\begin{equation}\label{m:eq:void-endpoint-norm}
 \left\|s\langle(\overline{J_v})_+\rangle_s\right\|_{2,k}
 \le C\sqrt{\frac{s}{\min(s,s_k)}+\delta_k s}.
\end{equation}
At $s=s_-$ this is at most $C(1+\sqrt{\delta_k s_k})$.
At $s=s_+$ the interpolation weight gives instead
\[
 C\frac{s_k}{T}\sqrt{\frac{T}{s_k}+\delta_k T}
 \le C\left(\sqrt{\frac{s_k}{T}}
       +\sqrt{\delta_k s_k}\sqrt{\frac{s_k}{T}}\right)
 \le C(1+\sqrt{\delta_k s_k}).
\]
The inner endpoint count term in \eqref{m:eq:void-endpoint-field} has
norm at most
\[
 C\sum_l2^{-l}(1+\sqrt{\delta_k2^lt})
 \le C(1+\sqrt{\delta_k t}).
\]
For the outer endpoint, $s_+/\max(s_l,s_+)\le1$; we retain its
contribution as $C(s_k/T)\bar L_{\mathrm{ann}}$.
Thus \eqref{m:eq:void-harnack} supplies a nonnegative, $X_k$-measurable
random variable $D_k$, the same for every electron in shell $k$,
such that
\begin{equation}\label{m:eq:void-final-field-bound}
 u_i\psi_k(x_i)\le D_k+C\frac{s_k}{T}\bar L_{\mathrm{ann}}-c\bar P_k,
 \qquad \|D_k\|_{2,k}\le C(1+\sqrt{\delta_k s_k}).
\end{equation}

\medskip\noindent
\emph{Summation over shells.}
Changing to the size-biased law and using
\eqref{m:eq:void-tilt-energy}, $s_k\ge1$, and
$m_k\le\E n_{\mathrm{neigh},k}$, gives
\begin{align}
 \E[n_kD_k]&\le m_k\E_kD_k
 \le C m_k(1+\sqrt{\delta_k s_k})\notag\\
 &\le C\left(m_k+
       \sqrt{m_k\E n_{\mathrm{neigh},k}/s_k}\right)
 \le C\E n_{\mathrm{neigh},k}.
 \label{m:eq:void-local-error-sum}
\end{align}
The neighbor annuli for supported $k$ lie inside $\{t\le|x|<T\}$ and
have bounded overlap, so these costs sum to $C\E L_{\mathrm{ann}}$.
For the outer term, $n_k$ is $X_k$-measurable and the tower property gives
\begin{align*}
 \sum_k\frac{s_k}{T}\E[n_k\bar L_{\mathrm{ann}}]
 &=\sum_k\frac{s_k}{T}\E[n_kL_{\mathrm{ann}}]\\
 &\le\|L_{\mathrm{ann}}\|_2
       \sum_k\frac{s_k}{T}\|n_k\|_2
 \le C\|L_{\mathrm{ann}}\|_2.
\end{align*}
Here the unconditioned norms use \eqref{m:eq:void-local-shell} with
$\delta=0$, and $\sum_k s_k/T\le1$.
Insert \eqref{m:eq:void-final-field-bound} into
\eqref{m:eq:void-weighted-field}.  The weights
$W_k=\sum_{i\in k}\chi(u_i)^2$ are $X_k$-measurable and satisfy
$0\le W_k\le n_k$.  Removing the bar on $P_k$ by the tower property,
we obtain
\begin{equation}\label{m:eq:void-prefix-products}
 \E\sum_kW_kP_k\le C(1+\|L_{\mathrm{ann}}\|_2).
\end{equation}

Let $I_{\mathrm{mid}}$ consist of the shells entirely contained in
$[64t,T/64]$, and set $L_{\mathrm{mid}}=\sum_{k\in I_{\mathrm{mid}}}n_k$.
On these shells $W_k=n_k$.  The sum on the left of
\eqref{m:eq:void-prefix-products} therefore includes all products
$n_kn_l$ with $k,l\in I_{\mathrm{mid}}$ and $l\le k-5$.
For every omitted close pair $|k-l|<5$, including $k=l$, use the tilt at
$k$ and \eqref{m:eq:void-local-shell} to obtain
\[
 \E n_kn_l\le m_k\E_kn_l
 \le C m_k(1+\sqrt{\delta_k s_l})
 \le C\E n_{\mathrm{neigh},k}.
\]
The last step uses $s_l\asymp s_k$ and the same calculation as
\eqref{m:eq:void-local-error-sum}.  Each $k$ has only a bounded number of
such neighbors.  Thus the close-pair costs sum to
$C\E L_{\mathrm{ann}}$, rather than a cost proportional to $J$.
Expanding $L_{\mathrm{mid}}^2$ now gives
\[
 \E L_{\mathrm{mid}}^2\le C(1+\|L_{\mathrm{ann}}\|_2).
\]
Only a fixed number of endpoint shells were omitted, so
$\|L_{\mathrm{ann}}-L_{\mathrm{mid}}\|_2\le C$ by
\eqref{m:eq:void-local-shell}.  Consequently
$X=\|L_{\mathrm{ann}}\|_2$ satisfies $X^2\le C(1+X)$, proving the lemma.
\end{proof}


\begin{lemma}[Exterior count]\label{m:lem:void-count}
There is a universal constant $C$ such that, for $0<s\le 1$,
\begin{equation}\label{m:eq:void-count}
 \E\#\{i:a_{x_i}\ge s\}\le CMs^{-3}.
\end{equation}
The constant is independent of the charges and the distances between
nuclei.
\end{lemma}

\begin{proof}
Write $\mathcal R=\{R_1,\ldots,R_M\}$ and
$d(x)=\operatorname{dist}(x,\mathcal R)$. We use
$d(x)\le L a_x$ from Lemma~\ref{m:lem:screen-scale}. For a Borel set $B$, write
$n(B)=\#\{i:x_i\in B\}$. The local count bound
\eqref{m:eq:screen-count}, applied with $\delta=0$, gives
\begin{equation}\label{m:eq:void-net-local-count}
 \E n(B(y,a_y))
 \le C\max\{a_y^{-3},1\}.
\end{equation}
We first prove that
\begin{equation}\label{m:eq:void-net-far-count}
 \E n(\{d\ge U\})\le CM
\end{equation}
for a sufficiently large universal constant $U$. A direct summation of
local bounds over all large distance scales would give an infinite
series. We instead group scales into a bounded number of annuli to which
Lemma~\ref{m:lem:void-shell} applies.

\paragraph{Nested sets of nuclear representatives.}
Fix
\[
 h=12,\qquad \varepsilon=2^{1-h}=\frac1{2048},\qquad
 D_0=1000,\qquad q=10.
\]
Thus
\[
 \varepsilon<\frac1{200},\qquad D_0-\varepsilon>400,
 \qquad D_0>3\varepsilon,\qquad 2^q>D_0.
\]
Choose an integer $j_0$ such that $2^{j_0}/L\ge1$ and
$2^{j_0-1}$ exceeds the lower threshold for the inner radius in
Lemma~\ref{m:lem:void-shell}. Set $U=2^{j_0}$.
Since the nuclei are distinct, there is an integer $l_0<j_0-h$ such that
$\mathcal R$ is $2^{l_0}$-separated. Put $P_{l_0}=\mathcal R$.
For each $l>l_0$, choose a maximal $2^l$-separated subset
$P_l\subset P_{l-1}$, where separated means that distinct points have
distance at least $2^l$.

For each $p\in P_{l-1}$ choose a parent $\pi_l(p)\in P_l$, keeping
$\pi_l(p)=p$ for $p\in P_l$. Maximality permits this choice with
$|p-\pi_l(p)|<2^l$ whenever $p$ is removed. Composing these maps gives
an ancestor $\alpha_l(R)\in P_l$ for every original nucleus $R$, and
\begin{equation}\label{m:eq:void-net-ancestor-offset}
 |R-\alpha_l(R)|<2^{l+1}.
\end{equation}
Indeed, the accumulated displacement is less than
$\sum_{k=l_0+1}^l2^k<2^{l+1}$; at $l=l_0$ it is zero.
Every representative is an original nucleus, and a removed nucleus never
reappears in the nested sets.

For $j\ge j_0$, put $u_j=2^j$. Choose a nearest nucleus $R(x)$
measurably, breaking ties by the least index. Assign a point of the dyad
$u_j\le d(x)<2u_j$ to
$c=\alpha_{j-h}(R(x))\in P_{j-h}$, and denote its assigned set by
$A_{j,c}$. These are Borel sets, disjoint at each scale, and
\begin{equation}\label{m:eq:void-net-assigned-set}
 |R(x)-c|<\varepsilon u_j,\qquad
 A_{j,c}\subset B(c,(2+\varepsilon)u_j).
\end{equation}

\paragraph{Only finitely many crowded pairs.}
Call $(j,c)$ crowded if there is a distinct
$w\in P_{j-h}$ with $|w-c|\le D_0u_j$; otherwise call it isolated.
We claim that the number of crowded pairs is at most $CM$.
For each crowded pair choose one witness $w$. The points $c$ and $w$
cannot both survive to $P_{j+q}$, since that set is
$2^{j+q}$-separated and $2^{j+q}>D_0u_j$. Let $r$ be the first level
at which either endpoint is removed, and charge $(j,c)$ to one such
removal event $(r,p)$, where $p\in P_{r-1}\setminus P_r$. Then
\begin{equation}\label{m:eq:void-net-charge-window}
 j-h<r\le j+q,
 \qquad\text{or equivalently}\qquad
 r-q\le j\le r+h-1.
\end{equation}

A fixed removal event can therefore receive charges from at most $q+h$
levels $j$. At any one of these levels, a center $c$ charging $(r,p)$
is either $p$ itself or has $p$ as its witness. In either case,
$c\in P_{j-h}\cap\overline B(p,D_0u_j)$. Such centers are
$2^{j-h}$-separated. The disjoint balls of radius $2^{j-h-1}$ around
them lie in the ball of radius $D_0u_j+2^{j-h-1}$ around $p$, so their
number is at most $(1+2D_0 2^h)^3$. This is a fixed constant. There are
at most $M-1$ removal events, since the nested sets eventually contain
one nucleus. Consequently
\begin{equation}\label{m:eq:void-net-crowded-number}
 \#\{(j,c):j\ge j_0,\ (j,c)\text{ is crowded}\}\le CM.
\end{equation}
Simultaneous removals cause no difficulty: the charge uses just one of
the two removed endpoints.

Each crowded assigned set has bounded expected particle count. To see
this, choose a maximal $u_j/(2L)$-separated subset of $A_{j,c}$.
By \eqref{m:eq:void-net-assigned-set}, a volume comparison bounds its cardinality
by $CL^3$. Its balls of radius $u_j/(2L)$ cover $A_{j,c}$.
At every center $y$ we have
$a_y\ge d(y)/L\ge u_j/L\ge1$, so these covering balls are contained
in $B(y,a_y)$. Equation~\eqref{m:eq:void-net-local-count} therefore gives
\begin{equation}\label{m:eq:void-net-assigned-count}
 \E n(A_{j,c})\le C.
\end{equation}
Together, \eqref{m:eq:void-net-crowded-number} and
\eqref{m:eq:void-net-assigned-count} show that all crowded assigned sets cost at
most $CM$ expected particles.

\paragraph{Isolated runs give empty annuli.}
For a fixed nucleus $c$, the levels $j\ge j_0$ at which
$c\in P_{j-h}$ form an initial integer interval, possibly empty or
infinite. Removing its crowded levels leaves consecutive runs of
isolated levels. The number of these runs is at most one plus the
number of crowded levels for $c$. Summing over $c$ and using
\eqref{m:eq:void-net-crowded-number}, the total number of isolated runs is at
most $CM$.

At an isolated level $j$ for $c$, every original nucleus $R$ satisfies
one of the two alternatives
\begin{equation}\label{m:eq:void-net-cluster-dichotomy}
 |R-c|<\varepsilon u_j
 \qquad\text{or}\qquad
 |R-c|>(D_0-\varepsilon)u_j.
\end{equation}
Indeed, if $\alpha_{j-h}(R)=c$, the first alternative follows from
\eqref{m:eq:void-net-ancestor-offset}. Otherwise isolation gives
$|\alpha_{j-h}(R)-c|>D_0u_j$, and the ancestor offset gives the
second alternative.

The alternative for a fixed nucleus cannot change between consecutive
isolated levels $j$ and $j+1$ for the same center $c$. A nucleus near
$c$ at level $j$ remains too close to satisfy the far alternative at
$j+1$. Conversely, a nucleus far at $j$ could be near at $j+1$ only if
its fixed distance from $c$ were both greater than
$(D_0-\varepsilon)u_j$ and less than $2\varepsilon u_j$, which is
impossible because $D_0>3\varepsilon$. Thus the set of near nuclei is
constant throughout each isolated run. This argument permits ancestor
maps to change and uses only the gap in
\eqref{m:eq:void-net-cluster-dichotomy}.

Consider first a finite run $j=a,\ldots,b$ and put
\[
 O=c,\qquad t=u_a/2,\qquad T=4u_b.
\]
The near nuclei have distance less than
$\varepsilon u_a<t/100$ from $c$. All other nuclei have distance
greater than $(D_0-\varepsilon)u_b>100T$. There is a nucleus at $c$,
and $T/t=2^{b-a+3}$ is a power of two. By the choice of $j_0$,
Lemma~\ref{m:lem:void-shell} applies. Every assigned point from this run
belongs to its annulus, since $c$ is itself a nucleus and
\[
 t\le u_a\le d(x)\le |x-c|<(2+\varepsilon)u_b<T.
\]
It follows that the union of the assigned sets over this run has
expected particle count at most $C$, independently of its length.
For an infinite run, apply the same argument to each finite initial
segment and then use monotone convergence for their increasing assigned
unions. There are at most $CM$ runs, so the total contribution of all
isolated assigned sets is at most $CM$. Combined with the crowded sets,
this proves \eqref{m:eq:void-net-far-count}.

\paragraph{The remaining distances.}
It remains to count particles with $a_x\ge s$ and $d(x)<U$.
For the part with $d(x)<s$, choose a maximal $s/2$-separated subset
of $\{d<s,\ a\ge s\}$. Its centers lie in
$\bigcup_{m=1}^M B(R_m,s)$. Disjoint balls of radius $s/4$ around
them show, by volume comparison, that there are at most $CM$ centers.
Their balls of radius $s/2$ cover the set and are contained in the
corresponding local balls $B(y,a_y)$, since $a_y\ge s$.
Equation~\eqref{m:eq:void-net-local-count} bounds this contribution by
$CMs^{-3}$.

For a distance dyad $u\le d(x)<2u$, choose a maximal
$u/(2L)$-separated subset. Its centers lie in
$\bigcup_m B(R_m,2u)$, so there are at most $CL^3M$ of them.
Their local balls cover the dyad, because $a_y\ge d(y)/L\ge u/L$.
Thus \eqref{m:eq:void-net-local-count} bounds its expected count by
\[
 CM\max\{(u/L)^{-3},1\}.
\]
Take $u=2^ks$ for $0\le k\le K$, where $K$ is the least nonnegative
integer such that $2^Ks\ge U$. These dyads cover $s\le d(x)<U$.
The sum of the inverse-cube terms is bounded by $Cs^{-3}$.
For the constant terms with $u<1$, use $1\le u^{-3}$; the number of
remaining dyads, with $1\le u\le2U$, is bounded by a universal constant.
Their total contribution is therefore at most $CMs^{-3}$.
Adding \eqref{m:eq:void-net-far-count} proves
\eqref{m:eq:void-count}.
\end{proof}
